\documentclass[11pt,a4paper]{article}
\input{../template/utm-lab}

\newcommand{\labnumber}{2}
\newcommand{\labtitle}{Meet the OS You Will Build}
\newcommand{\studentname}{Daniel Nicorici}
\newcommand{\studentgroup}{FAF-243}

\begin{document}

\utmtitlepage
\tableofcontents
\newpage

\plainsection{Purpose of the work}

The purpose of this lab is to build xv6, a small but real Unix written in about 9\,000 lines
of C, from source; to boot it on an emulated RISC-V machine in QEMU; to use it as the tiny Unix
it is; to read both sides of a system call (the program in \code{user/} that asks and the
kernel code in \code{kernel/} that answers); and to add a first program of my own,
\code{sleep}, that makes a real system call into the kernel I compiled.

Environment: xv6 source \code{mit-pdos/xv6-riscv} at commit \code{06aad25} (2026-09-26),
cross-compiled with \code{riscv64-linux-gnu-gcc} 16.2.1 and run on QEMU 11.1.1
(\code{qemu-system-riscv64}), on my own laptop running Arch Linux.

\section{Part 1. Build it and boot it}

\begin{lstlisting}
$ cd ~/xv6-riscv
$ make qemu
qemu-system-riscv64 -machine virt -bios none -kernel kernel/kernel -m 128M -smp 3 -nographic -global virtio-mmio.force-legacy=false -drive file=fs.img,if=none,format=raw,id=x0 -device virtio-blk-device,drive=x0,bus=virtio-mmio-bus.0

xv6 kernel is booting

hart 2 starting
hart 1 starting
init: starting sh
$
\end{lstlisting}

\code{make qemu} compiles the kernel into \code{kernel/kernel}, compiles every user program
listed in \code{UPROGS}, packs \code{README} and those programs into the disk image
\code{fs.img}, and starts QEMU with that kernel and disk. The emulated machine is the RISC-V
\code{virt} board with 128\,MB of RAM and 3 harts (hardware threads, i.e. CPU cores): hart 0
boots the kernel, harts 1 and 2 report that they are starting, and then the first process,
\code{init}, starts the shell. QEMU is left with \code{Ctrl-A} followed by \code{X}.

\section{Part 2. Use xv6 as the Unix it is}

\begin{lstlisting}
$ ls
.              1 1 1024
..             1 1 1024
README         2 2 2441
cat            2 3 35912
echo           2 4 34808
forktest       2 5 17704
grep           2 6 43336
init           2 7 35152
kill           2 8 34704
ln             2 9 34504
ls             2 10 41992
mkdir          2 11 34760
rm             2 12 34736
sh             2 13 57120
stressfs       2 14 35584
usertests      2 15 201568
grind          2 16 50384
wc             2 17 36880
zombie         2 18 34056
logstress      2 19 36728
forphan        2 20 35520
dorphan        2 21 34976
sync           2 22 34152
sleep          2 23 34480
console        3 24 0
$ echo hello xv6
hello xv6
$ ls | grep c
cat            2 3 35912
echo           2 4 34808
wc             2 17 36880
sync           2 22 34152
console        3 24 0
$ wc README
48 336 2441 README
\end{lstlisting}

The columns printed by xv6's \code{ls} are: name, type (1 = directory, 2 = file, 3 = device),
inode number and size in bytes. \code{console} has type 3: the terminal is a device file, just
like \code{/dev/tty} on Linux.

\textbf{Observations:}
\begin{enumerate}
  \item \textbf{Three programs xv6 ships with.} \code{ls}, \code{cat} and \code{grep}; the
  image also contains \code{echo}, \code{wc}, \code{sh}, \code{mkdir}, \code{rm},
  \code{kill}, \code{ln} and the test programs \code{usertests} and \code{forktest}.
  \item \textbf{What a pipe needs.} First, \emph{processes}: the shell uses \code{fork} and
  \code{exec} to create two processes, one running \code{ls} and one running \code{grep}, which
  run at the same time. Second, \emph{inter-process communication}: the \code{pipe} system call
  creates a kernel buffer with a write end and a read end, and the shell uses \code{close} and
  \code{dup} to connect the standard output of \code{ls} to the write end and the standard
  input of \code{grep} to the read end (\code{user/sh.c}, lines 103--113).
  \item \textbf{xv6 shell vs.\ Linux shell.} For simple commands the xv6 shell behaves like
  bash (programs, pipes, \code{>}/\code{<} redirection, \code{;} and \code{\&}), but it is much
  more minimal: no history, no Tab completion, no variables, no \code{*} globbing, and the
  prompt is just \code{\$}.
\end{enumerate}

\section{Part 3. Read the source}

\subsection{System calls in user/cat.c}

\begin{lstlisting}[style=ccode]
#include "kernel/types.h"
#include "kernel/fcntl.h"
#include "user/user.h"

char buf[512];

void
cat(int fd)
{
  int n;

  while ((n = read(fd, buf, sizeof(buf))) > 0) {
    if (write(1, buf, n) != n) {
      fprintf(2, "cat: write error\n");
      exit(1);
    }
  }
  if (n < 0) {
    fprintf(2, "cat: read error\n");
    exit(1);
  }
}
\end{lstlisting}

\subsection{Where the kernel implements read and write}

\begin{lstlisting}
$ grep -n "sys_read\|sys_write" kernel/sysfile.c
69:sys_read(void)
83:sys_write(void)
\end{lstlisting}

\begin{lstlisting}[style=ccode,firstnumber=68]
uint64
sys_read(void)
{
  struct file *f;
  int n;
  uint64 p;

  argaddr(1, &p);
  argint(2, &n);
  if (argfd(0, 0, &f) < 0)
    return -1;
  return fileread(f, p, n);
}
\end{lstlisting}

\subsection{How a call crosses from user/ into kernel/}

On the user side, \code{read} is not C code at all. The script \code{user/usys.pl} generates
a three-instruction stub for every system call in \code{user/usys.S}: it loads the call
number into register \code{a7} and executes \code{ecall}, which traps into the kernel.

\begin{lstlisting}
$ grep -A3 "^read:" user/usys.S
read:
 li a7, SYS_read
 ecall
 ret
\end{lstlisting}

In the kernel, \code{syscall()} reads that number back from the saved registers and calls the
matching function from its table (\code{[SYS\_read] = sys\_read} on line 115, with
\code{SYS\_read} = 5 defined in \code{kernel/syscall.h}). The return value is stored in
\code{a0}, which is where the user program finds it.

\begin{lstlisting}[style=ccode,firstnumber=136]
void
syscall(void)
{
  int num;
  struct proc *p = myproc();

  num = p->trapframe->a7;
  if (num > 0 && num < NELEM(syscalls) && syscalls[num]) {
    // Use num to lookup the system call function for num, call it,
    // and store its return value in p->trapframe->a0
    p->trapframe->a0 = syscalls[num]();
  } else {
    printk("%d %s: unknown sys call %d\n", p->pid, p->name, num);
    p->trapframe->a0 = -1;
  }
}
\end{lstlisting}

\textbf{Observations:}
\begin{enumerate}
  \item \textbf{System calls used by \code{user/cat.c}.} \code{open(path, O\_RDONLY)} asks the
  kernel to find a file by name and return a file descriptor; \code{read(fd, buf, 512)} asks
  it to copy up to 512 bytes from that file into \code{buf} (it returns the number of bytes,
  0 at end of file, $-1$ on error); \code{write(1, buf, n)} asks it to copy \code{n} bytes to
  file descriptor 1, the console; \code{close(fd)} releases the descriptor; \code{exit(status)}
  ends the process and passes the status to the parent's \code{wait}. \code{fprintf} is not a
  system call: it is a library function in \code{user/printf.c} that formats the text and then
  calls \code{write} on descriptor 2.
  \item \textbf{Where \code{sys\_read} is implemented.} In \code{kernel/sysfile.c}, line 69
  (with its return type \code{uint64} on line 68); \code{sys\_write} follows on line 83.
  \item \textbf{\code{kernel/} vs.\ \code{user/}.} The code in \code{kernel/} runs in
  privileged (supervisor) mode with access to all memory and hardware and provides the
  services; the code in \code{user/} runs as ordinary, isolated programs that can only obtain
  those services by asking the kernel through system calls.
\end{enumerate}

\section{Part 4. Write your first xv6 program}

\subsection{Implementation}

In August 2025 the xv6 authors renamed the user system call \code{sleep()} to
\code{pause()} (upstream commit \code{f5dea58}, which also renamed kernel functions such as
\code{fork} to \code{kfork} and \code{exit} to \code{kexit}). On the current tree
\code{user/user.h} declares only \code{int pause(int);}, so the handout's version, which calls
\code{sleep()}, does not build. My program therefore calls \code{pause()}:

\begin{lstlisting}[style=ccode]
#include "kernel/types.h"
#include "user/user.h"

int
main(int argc, char *argv[])
{
  if(argc != 2){
    fprintf(2, "usage: sleep <ticks>\n");
    exit(1);
  }
  pause(atoi(argv[1]));
  exit(0);
}
\end{lstlisting}

\code{argv[0]} is the program name and \code{argv[1]} the number of ticks, so any other
argument count is a usage error: the message goes to file descriptor 2 (standard error) and
\code{exit(1)} reports the failure to the parent shell.

\subsection{Adding the program to the build}

\begin{lstlisting}
$ git diff Makefile
diff --git a/Makefile b/Makefile
index af1e7a5..089e215 100644
--- a/Makefile
+++ b/Makefile
@@ -150,6 +150,7 @@ UPROGS=\
 	$U/_forphan\
 	$U/_dorphan\
 	$U/_sync\
+	$U/_sleep\

 fs.img: mkfs/mkfs README $(UPROGS)
 	mkfs/mkfs fs.img README $(UPROGS)
\end{lstlisting}

User programs are built as \code{user/\_name}; \code{mkfs} strips the leading underscore when
it copies them into \code{fs.img} (\code{mkfs/mkfs.c}, line 147), so inside xv6 the program is
simply \code{sleep}. Without the \code{UPROGS} line the program is never compiled or copied, and
the shell answers \code{exec sleep failed}.

\subsection{Testing}

\begin{lstlisting}
$ sleep 10
$ sleep
usage: sleep <ticks>
$ sleep abc 2
usage: sleep <ticks>
\end{lstlisting}

To check that the program really waits rather than returning at once, I timed it from the
host: after typing \code{sleep 50; echo done} the prompt line appeared at 41.46\,s and
\code{done} at 46.42\,s, a pause of 4.96\,s. One xv6 tick is about 0.1\,s (the kernel arms
the next timer interrupt 1\,000\,000 cycles ahead, \code{kernel/trap.c}, line 179), so 50 ticks
$\approx$ 5\,s, as expected.

\plainsection{Conclusions}

I built xv6 from source, booted it in QEMU and used it as a small Unix with the same basic tools
and pipes as Linux. Reading the source showed the split at the heart of every operating system:
programs in \code{user/} only ask, and the kernel in \code{kernel/} does the privileged work.

My \code{sleep} program crosses that boundary for real. When it calls \code{pause(10)}, the
argument goes into register \code{a0}, and the stub in \code{user/usys.S} loads the call number
13 into \code{a7} and executes \code{ecall}. The CPU switches to supervisor mode,
\code{uservec} saves the user registers, \code{usertrap()} recognises a system call, and
\code{syscall()} uses \code{a7} to call \code{sys\_pause}, which waits for ten timer ticks. The
result travels back in \code{a0}, the CPU returns to user mode, and the program exits. Every
service a program gets from the OS, from reading a file to sleeping, follows this same path.

\end{document}
